\documentclass[12pt,a4paper]{article}
\usepackage[a4paper,left=2.5cm,right=2cm,top=2cm,bottom=2cm]{geometry}
\usepackage{fontspec}
\setmainfont{DejaVu Serif}
\setmonofont{DejaVu Sans Mono}
\usepackage{setspace}
\onehalfspacing
\usepackage{enumitem}
\usepackage{booktabs}
\usepackage{tabularx}
\usepackage{array}
\usepackage{xcolor}
\usepackage{listings}
\usepackage{hyperref}
\hypersetup{hidelinks}
\usepackage{titlesec}
\titleformat{\section}{\normalfont\Large\bfseries}{\thesection.}{0.5em}{}
\titleformat{\subsection}{\normalfont\large\bfseries}{\thesubsection.}{0.5em}{}
\setlength{\parindent}{1.25cm}
\setlength{\parskip}{0pt}
\renewcommand{\contentsname}{Contents}
\definecolor{codebg}{RGB}{246,246,246}
\definecolor{codeframe}{RGB}{215,215,215}
\lstset{basicstyle=\ttfamily\small,backgroundcolor=\color{codebg},frame=single,
  rulecolor=\color{codeframe},breaklines=true,columns=fullflexible,keepspaces=true,
  showstringspaces=false,aboveskip=0.8em,belowskip=0.5em}
\newcommand{\figcap}[1]{\begin{center}\small\textit{#1}\end{center}\vspace{0.4em}}

\begin{document}
\begin{titlepage}
\begin{center}
{\large Ministerul Educației și Cercetării al Republicii Moldova}\\[0.3cm]
{\large Universitatea Tehnică a Moldovei}\\[0.3cm]
{\large Facultatea Calculatoare, Informatică și Microelectronică}\\[0.3cm]
{\large Departamentul Ingineria Software și Automatică}\\[0.3cm]
{\large Program: Software Engineering}\\[2.0cm]
{\Large \textbf{Operating Systems}}\\[0.8cm]
{\LARGE \textbf{Laboratory Work No. 2}}\\[0.5cm]
{\Large \textbf{Meet the OS You Will Build}}\\[0.35cm]
{\large Track B: Build Your Own Small OS (xv6)}
\vfill
\begin{flushright}
\begin{tabular}{ll}
\textbf{Done by:} & st. gr. FAF-241, Ermicev Sofia \\
\\
\textbf{Checked by:} & Patricia Reițman
\end{tabular}
\end{flushright}
\vfill
{\large Chișinău -- 2026}
\end{center}
\end{titlepage}
\tableofcontents
\newpage

\section{Purpose of the Work}
The purpose of this laboratory work is to build and boot xv6, use it as a small Unix-like operating system, and connect observable shell behavior with the source code in the \texttt{user/} and \texttt{kernel/} directories.

\section{Part 1. Building and Booting xv6}
The xv6 source tree was built with the RISC-V toolchain and started in QEMU. The actual boot reached the shell successfully.
\begin{lstlisting}
xv6 kernel is booting

init: starting sh
$
\end{lstlisting}
\figcap{Figure 1.1 -- Actual xv6 boot excerpt showing that the shell started successfully}

The build compiled kernel and user programs, created the file-system image with \texttt{mkfs}, and linked the kernel. A linker warning about an RWX segment was observed, but the kernel booted and the shell was usable.

\section{Part 2. Use xv6 as the Unix It Is}
The xv6 image contains ordinary user programs such as \texttt{ls}, \texttt{cat}, and \texttt{echo}. Basic command execution and pipes behave similarly to Unix/Linux even though the shell is much smaller.

\subsection{Observe answers}
\begin{enumerate}[label=\arabic*.]
\item Three programs included with xv6 are \texttt{ls}, \texttt{cat}, and \texttt{echo}.
\item For a pipeline such as \texttt{ls | grep c} to work, the OS needs processes and inter-process communication. The shell starts separate processes and connects the output of one process to the input of another through a pipe.
\item The xv6 shell is much simpler than Bash, but basic Unix mechanisms such as program execution, standard input/output, and pipes work according to the same core model.
\end{enumerate}

\section{Part 3. Reading the Source}
The source of \texttt{user/cat.c} and the kernel system-call implementation were inspected.
\begin{lstlisting}
sys_read  -> kernel/sysfile.c, line 69
sys_write -> kernel/sysfile.c, line 83
\end{lstlisting}
\figcap{Figure 3.1 -- Locations of the read and write system-call handlers found in the xv6 source}

\subsection{Observe answers}
\begin{enumerate}[label=\arabic*.]
\item \texttt{user/cat.c} uses \texttt{read()} to obtain bytes from a file descriptor, \texttt{write()} to send bytes to a file descriptor, \texttt{open()} to open a file, \texttt{close()} to release a descriptor, and \texttt{exit()} to terminate the process.
\item \texttt{sys\_read} is implemented in \texttt{kernel/sysfile.c} at line 69, and \texttt{sys\_write} starts at line 83 in the inspected tree.
\item Code under \texttt{kernel/} runs as privileged operating-system code and manages resources. Code under \texttt{user/} consists of ordinary programs that request kernel services through system calls.
\end{enumerate}

\section{Conclusion}
This work confirmed that xv6 is a small but real Unix-like operating system. Building and booting it exposed the complete path from source code to a running shell, while the source inspection connected user-level calls such as \texttt{read()} and \texttt{write()} with privileged kernel implementations.

\end{document}
